%! Solving Polynomial Equations by Factoring — Unit 1, Lesson 1.6 — Warm-Up (Answer Key)
\documentclass[10pt]{article}
\usepackage{atda-article}
\usepackage{atda-key}   % pulls in -boxes; do NOT also load -boxes
\newcommand{\TallMath}[1]{$\displaystyle #1\rule[-1.4em]{0pt}{3.2em}$}
\setlength{\workrowsep}{6pt}   % handwriting room; moves blank and key together

\begin{document}

\pageheader{Unit 1, Lesson 1.6}{Warm-Up --- Answer Key}
% NAMESTRIP: no \namedateperiod here — mirrors the blank. Teacher-only prose goes
% in the lesson plan as \begin{teachernote}[Warm-Up], never in a key.

\vspace{0.15in}

\begin{notesbox}{Four Moves Today's Lesson Runs On}
\small
Item 4 uses your answers from items 2 and 3 --- do them in order. Show your work on all four.

\begin{enumerate}[leftmargin=1.6em, itemsep=6pt, topsep=4pt]

\item \textbf{Spiral --- Lessons 1.4 and 1.5.} Factor: \quad (a) $x^2-7x+12$ \quad
(b) $x^2-49$
\begin{work}
  x^2-7x+12 &= (x-3)(x-4) \\
  x^2-49 &= (x+7)(x-7)
\end{work}

\item \textbf{Spiral --- Lesson 1.1.} Evaluate $(x-3)(x+5)$ at $x=3$, at $x=-5$, and at $x=1$.
\begin{work}
  (3-3)(3+5) &= (0)(8) = 0 \\
  (-5-3)(-5+5) &= (-8)(0) = 0 \\
  (1-3)(1+5) &= (-2)(6) = -12
\end{work}

\item \textbf{Number sense --- the whole lesson is hiding here.} Suppose $a$ and $b$ are numbers.
(a) If $ab=0$ and $a=4$, what must $b$ be? \ (b) If $ab=12$ and $a=4$, what must $b$ be? \
(c) If all you are told is $ab=12$, can you name $a$? Which of the two products, $0$ or $12$,
forces a factor to be something in particular?
\begin{work}
  4b &= 0 \quad\Rightarrow\quad b = 0 \\
  4b &= 12 \quad\Rightarrow\quad b = 3 \\
  ab &= 12 \quad\Rightarrow\quad a \text{ could be } 1, \, 2, \, 3, \, 12, \, \tfrac12, \dots
        \text{ --- no conclusion} \\
  ab &= 0 \quad\Rightarrow\quad a = 0 \text{ or } b = 0 \text{ --- only the product } 0 \text{ forces it}
\end{work}

\item \textbf{Discovery.} Item 2 found two values of $x$ making $(x-3)(x+5)$ equal to zero. Using
item 3, explain how you could have found both of them \emph{without} testing numbers --- and say why
there cannot be a third.
\begin{work}
  x-3 &= 0 \quad\Rightarrow\quad x = 3 \\
  x+5 &= 0 \quad\Rightarrow\quad x = -5 \\
  \text{any other } x: \; & \text{both factors are nonzero, and a product of} \\
                         & \text{two nonzero numbers is never zero}
\end{work}

\end{enumerate}
\end{notesbox}

\end{document}
